\chapter{Reproducibilidad}
\label{ap:reproducibilidad}

Este apéndice recoge lo necesario para reproducir cualquier cifra de la memoria. No es
documentación de instalación: es la respuesta a la pregunta de si los resultados
presentados pueden volver a obtenerse y con qué grado de fidelidad.

\section{Entorno de ejecución}

Los experimentos se ejecutaron sobre una tarjeta gráfica AMD con la pila de cómputo ROCm,
en un sistema Linux. Esta elección tiene una consecuencia que conviene declarar: parte del
software habitual del campo asume tarjetas de otro fabricante, y algunas implementaciones
optimizadas para inferencia no estaban disponibles. La comparativa se realizó por tanto
sobre la implementación de referencia de la biblioteca \emph{Transformers}
\cite{wolf2020transformers}, más lenta que las alternativas especializadas pero disponible
por igual en ambas plataformas, lo que evita que la elección de hardware se confunda con una
diferencia de resultados.

El intérprete de Python se creó reutilizando los paquetes del sistema de forma deliberada,
porque reinstalar la biblioteca de cálculo dentro del entorno aislado rompe el soporte de
la tarjeta gráfica. Es un detalle menor que causa horas de depuración si se desconoce.

\section{Trazabilidad de los resultados}

El procedimiento actual hace que cada ejecución produzca un fichero de resultados que
registra, junto a las métricas:

\begin{itemize}
  \item El identificador del commit del código que la produjo.
  \item Si el árbol de trabajo tenía cambios sin confirmar en ese momento, lo que indica
        que el commit registrado no basta para reconstruir el estado exacto.
  \item El resumen criptográfico del manifiesto de corpus empleado, junto con su número de
        fragmentos, y el de la submuestra efectiva cuando el experimento evalúa solo una
        parte de él.
  \item La configuración completa: modelo, estrategia de decodificación, semilla,
        parámetros del experimento y versiones de las bibliotecas relevantes.
\end{itemize}

El tercer punto se incorporó tras un incidente. Al ampliar un corpus se sobrescribió su
manifiesto, y los resultados anteriores quedaron apuntando a un fichero cuyo contenido ya
era distinto sin que nada lo advirtiera. El resumen criptográfico permite detectar esa
discrepancia; posteriormente se añadió también la posibilidad de que varias muestras de la
misma fuente coexistan, que la previene.

\subsection*{Hasta dónde llega, en la práctica}

Decir «cada ejecución registra» describe el procedimiento, no el archivo: el mecanismo
nació de aquel incidente, de modo que alcanza en distinto grado a lo ejecutado antes y
después. Conviene dar las cifras, porque presentar la trazabilidad como una propiedad
uniforme sería justamente el tipo de afirmación de la que este trabajo dedica un capítulo a
desconfiar. Sobre los cincuenta y dos ficheros de resultados que sostienen esta memoria:

\begin{itemize}
  \item Todos registran su configuración completa, que es lo que permite repetir la
        ejecución, y todos registran un identificador de commit.
  \item En cuarenta y cuatro ese identificador vale \texttt{sin-git}, porque se produjeron
        antes de que el trabajo estuviera bajo control de versiones. Son los más antiguos,
        y entre ellos está el barrido de línea base.
  \item Once registran el resumen criptográfico del manifiesto, que es la parte que se
        añadió tras el incidente. En los once, recalcularlo hoy devuelve el valor que el
        fichero declara.
  \item Uno solo se produjo con el árbol de trabajo limpio. El resto de los que lo declaran
        tenían cambios sin confirmar, cosa que el propio bloque advierte y que significa
        que el commit registrado no basta para reconstruir el estado exacto.
\end{itemize}

El programa \texttt{tools/auditar\_resultados.py} recorre el archivo entero y falla si algo
no cuadra. Aplica dos comprobaciones, según lo que cada fichero permita.

La directa consiste en \emph{recalcular} el resumen criptográfico del manifiesto y
compararlo con el que el resultado declara. Merece subrayarse por lo que tiene de lección
repetida: guardar el resumen no detecta nada por sí solo, y durante un tiempo nadie lo
recalculó, de modo que la salvaguarda contra el incidente era exactamente la misma clase de
fallo que este trabajo documenta en otros sitios, un dato declarado que ninguna
comprobación consulta. Cuando el experimento evalúa solo una parte del manifiesto se
compara además el resumen de esa submuestra, porque el del fichero completo identifica
material que no se ha medido.

La indirecta cubre a los resultados anteriores al mecanismo, que no llevan resumen: el
número de fragmentos y de palabras de referencia que declara cada fichero tiene que
coincidir con el que produce el manifiesto de hoy. Hoy señala cinco ficheros, y son
precisamente los del incidente: mediciones sobre cuarenta fragmentos de un corpus que
después se amplió a mil doscientos. Ese es el motivo de que la
tabla~\ref{tab:baseline-corpus} marque esa fila, y de que las comparaciones entre modelos
deban hacerse sobre las tablas por corpus.

Dos categorías quedan fuera de la comprobación indirecta sin que eso sea un hueco. Los
barridos de ingeniería vuelven a trocear el audio en ventanas o segmentos, así que no
declaran un recuento de fragmentos con el que contrastar; llevan resumen criptográfico, que
es la comprobación fuerte. Y los experimentos que emplean glosario reservan una fracción de
los fragmentos para extraerlo y evalúan sobre el resto, de modo que su recuento queda por
debajo del manifiesto a propósito: el programa descuenta esa reserva, porque denunciar como
deriva la propia salvaguarda contra la fuga de información habría enterrado los cinco casos
reales bajo cuatro falsos.

Merece una advertencia de uso, porque costó tiempo: el recuento de palabras que guardan los
resultados es el de las referencias \emph{ya normalizadas}, y el normalizador convierte los
signos en espacios, de modo que un corpus puede declarar más palabras de las que tiene en
crudo. Compararlo contra el texto sin normalizar produce discrepancias que no existen.

\subsection*{Qué se versiona y qué no}

Los manifiestos de corpus, los programas y las tablas y figuras generadas están en el
repositorio. No lo están, deliberadamente, tres cosas: el audio, porque su licencia no
permite redistribuirlo; los ficheros de resultados en crudo, porque son la salida de los
programas y no su entrada; y el adaptador entrenado, por tamaño.

La consecuencia hay que decirla sin adornos: \textbf{regenerar cualquier cifra con un solo
comando es cierto sobre una instalación que ya tenga el corpus descargado}, que es el
supuesto en el que se ha trabajado. Un tercero que parta del repositorio necesita antes
obtener el audio de sus fuentes originales mediante los programas de descarga incluidos,
disponer de una tarjeta gráfica, y reentrenar el adaptador para reproducir el resultado del
ajuste fino. Los manifiestos hacen esa reconstrucción posible y verificable, que es su
función; no la hacen inmediata.

\section{Identidad de un resultado}

Los ficheros de resultados se nombran según el modelo, el corpus y la estrategia de
decodificación empleados. Esta última forma parte del nombre porque se comprobó que dos
ejecuciones idénticas salvo por ese parámetro producían tasas de error que diferían en más
de doce puntos. Un resultado identificado únicamente por modelo y corpus habría permitido
comparar cifras no comparables.

La regla general que se derivó: cualquier eje de variación que se introduzca debe
incorporarse al nombre del resultado, no solo a su contenido.

\section{Generación de tablas y figuras}

Ninguna cifra de esta memoria se ha transcrito manualmente. Las tablas y figuras se generan
mediante programas que leen los ficheros de resultados y producen directamente el código
LaTeX o el gráfico correspondiente, que la memoria incorpora por inclusión.

La consecuencia práctica es que regenerar los informes tras una nueva ejecución actualiza
la memoria sin intervención. La consecuencia metodológica es más relevante: una cifra
copiada a mano deja de estar vinculada a la ejecución que la produjo en el momento en que
se copia, y no hay forma de detectar la divergencia posterior.

\section{Verificaciones previas}

Tres comprobaciones se ejecutan antes de dar por válido un resultado, y las tres surgieron
de errores que no producían ningún fallo visible.

\textbf{Verificación de la tarjeta gráfica.} No basta con comprobar que el sistema la
detecta: se verifica que calcula correctamente y que completa un ciclo de entrenamiento con
descenso efectivo de la función de pérdida.

\textbf{Verificación del entrenamiento.} Antes de evaluar un modelo ajustado se comprueba
que la función de pérdida se encuentra en el rango esperado y que desciende entre épocas.
Un ajuste con las etiquetas mal alineadas termina sin errores, guarda su resultado y
produce métricas evaluables; solo el valor absoluto de la pérdida lo delata.

\textbf{Verificación del procesamiento por lotes.} Agrupar fragmentos acelera
sustancialmente los barridos, pero solo es admisible si no altera las transcripciones. Se
comprueba comparándolas con las obtenidas de una en una. Una optimización que cambia los
resultados no es una optimización: es un factor de confusión.

\section{Estructura del repositorio}

El trabajo separa estrictamente el material de investigación del de la aplicación; ambos
no comparten código y se comunican mediante una interfaz documentada. Cada experimento
reside en su propia carpeta, con el programa que lo ejecuta, una descripción de su
hipótesis y su diseño, y sus resultados.

Los manifiestos de corpus se versionan; el audio no. Un manifiesto contiene las rutas, las
transcripciones de referencia, las duraciones y los metadatos disponibles, de modo que el
material puede reconstruirse desde su fuente original sin almacenar los ficheros de sonido.
